\documentclass[10pt,a4paper]{article}

\usepackage[a4paper,top=18mm,bottom=19mm,left=21mm,right=21mm]{geometry}
\usepackage{fontspec}
\usepackage{amsmath,amssymb}
\usepackage{booktabs,tabularx,array}
\usepackage{xcolor}
\usepackage{hyperref}
\usepackage{enumitem}
\usepackage{caption}

\setmainfont{Times New Roman}
\setsansfont{Helvetica Neue}
\setmonofont{Menlo}[Scale=0.82]
\definecolor{navy}{HTML}{17365D}
\definecolor{blue}{HTML}{1F5A94}
\definecolor{lightblue}{HTML}{EAF2F8}
\hypersetup{
  colorlinks=true,
  linkcolor=navy,
  citecolor=navy,
  urlcolor=blue,
  pdftitle={Báo cáo tuần: MPPI bám global path và điều khiển goal từ RViz},
  pdfauthor={Nguyễn Thành Lâm}
}
\setlength{\parindent}{1.1em}
\setlength{\parskip}{0.18em}
\setlist{nosep,leftmargin=1.45em}
\captionsetup{font=small,labelfont=bf}
\newcommand{\code}[1]{\texttt{#1}}
\newcolumntype{Y}{>{\raggedright\arraybackslash}X}
\newcolumntype{L}[1]{>{\raggedright\arraybackslash}p{#1}}

\title{\vspace{-2.2em}\textbf{Báo cáo tiến độ tuần\\
MPPI bám global path và điều khiển goal từ RViz}}
\author{Nguyễn Thành Lâm\\[-0.1em]
\small Companion-computer local planning cho ArduPilot}
\date{13 tháng 09 năm 2026}

\begin{document}
\maketitle
\vspace{-1.7em}

\begin{center}
\fcolorbox{navy}{lightblue}{\parbox{0.91\textwidth}{\centering
\textbf{Video demo:} \href{https://drive.google.com/file/d/1I1VxN6-6Y3Gt8kxY8ieNvlvMMkvvLxLJ/view?usp=sharing}{Google Drive}\quad
\textbf{Mã nguồn:} \href{https://github.com/thanhlamauto/ardu_plugin_gazebo}{GitHub -- thanhlamauto/ardu\_plugin\_gazebo}}}
\end{center}

\begin{abstract}
Tuần này hệ thống local planning trên companion computer được nâng từ MPPI
điểm-đến-điểm lên MPPI bám global path theo thời gian, đồng thời bổ sung giao
diện chọn goal trực tiếp trong RViz. Planner nhận odometry và point cloud LiDAR
3D, sinh velocity/yaw-rate setpoint rồi gửi ArduPilot qua MAVLink. Ba tuyến
right-angle, narrow-gate và slalom đã đạt terminal trong Gazebo/SITL với một
seed mỗi map. Kết quả cho thấy pipeline hoạt động, nhưng chưa chứng minh độ bền
đa seed, rigid-body PA-MPPI hay an toàn phần cứng.
\end{abstract}

\noindent\textbf{Từ khóa:} MPPI, global-path tracking, RViz, MAVLink,
ArduPilot, Gazebo, LiDAR 3D.

\section{Mục tiêu và kết quả chính}

Mục tiêu kỹ thuật là giữ perception và local optimization ở companion
computer, còn ArduPilot đảm nhiệm stabilization và tracking setpoint. Đây là
lựa chọn kiến trúc của project; không hàm ý stack tránh vật cản của ArduPilot
không sử dụng được. Guided mode và MAVLink cung cấp interface velocity target
phù hợp cho cách phân tầng này \cite{ardupilotGuidedDocs2026,mavlinkCommon2026}.

Các kết quả chính trong tuần:
\begin{enumerate}
  \item thêm cost bám global path theo thời gian dựa trên cấu trúc reference
  tracking của Minařík et al. \cite{minarik2024agilemppi};
  \item mô phỏng cùng command conditioner trong rollout và tại MAVLink output;
  \item thêm predictive brake, timeout/stale-data gate và terminal handling;
  \item hiển thị nominal trajectory cùng top-$K$ sampled trajectories trong RViz;
  \item nhận \code{/goal\_pose} từ công cụ \code{2D Goal Pose}, cho phép thay
  goal mà không dùng thao tác \code{Fly To};
  \item tạo profile demo mượt riêng, không sửa ngược profile dùng cho audit.
\end{enumerate}

\begin{table}[h]
\centering\small
\caption{Trạng thái bằng chứng tại thời điểm báo cáo.}
\begin{tabularx}{\textwidth}{L{0.30\textwidth}Y}
\toprule
\textbf{Trạng thái} & \textbf{Phạm vi}\\
\midrule
Verified offline & 41/41 unit tests; deterministic point-mass smoke test.\\
Verified Gazebo/SITL, one seed & Velocity-level MPPI trên ba fixed routes.\\
Implemented, chưa live-verified & RViz-click flight hoàn chỉnh; profile demo mượt mới.\\
Pending & Rigid-body hover/body-rate, PA-MPPI unknown-space ablation, edge computer và hardware.\\
\bottomrule
\end{tabularx}
\end{table}

\section{Kiến trúc và luồng thực thi}

\begin{center}
\fcolorbox{navy}{lightblue}{\parbox{0.90\textwidth}{\centering
LiDAR 3D + Gazebo odometry
$\rightarrow$ đổi frame/downsample
$\rightarrow$ MPPI hoặc PA-MPPI
$\rightarrow$ MAVLink velocity setpoint
$\rightarrow$ ArduPilot GUIDED
$\rightarrow$ motors/Gazebo}}
\end{center}

Baseline dùng trạng thái vật lý
$x=[p_x,p_y,p_z,v_x,v_y,v_z,\psi]$ và lệnh
$u=[v_x^{cmd},v_y^{cmd},v_z^{cmd},\dot\psi^{cmd}]$. Rollout còn lưu lệnh đã
qua conditioner ở bước trước để mô hình dự đoán đúng giới hạn thay đổi lệnh.
Mô hình này cố ý giảm bậc: attitude và motor dynamics vẫn do ArduPilot xử lý.

Các module chính gồm:
\begin{itemize}
  \item \code{mppi\_controller.py}: dynamics, sampling, cost và trajectory;
  \item \code{mppi\_local\_planner\_node.py}: vòng điều khiển, safety gates,
  waypoint, RViz và logging;
  \item \code{mavlink\_interface.py}: chuyển ENU/NED và gửi setpoint;
  \item \code{scripts/mppi\_velocity\_avoidance.py}: CLI và nạp YAML;
  \item \code{occupancy\_grid.py}, \code{pa\_mppi\_controller.py}: nhánh
  perception-aware experimental.
\end{itemize}

RViz chỉ gửi \code{PoseStamped}; node lấy $x,y$ từ click và giữ độ cao hiện tại
hoặc dùng độ cao cấu hình. Mỗi click reset route, terminal state, conditioner
và MPPI warm start. Arm, takeoff và emergency mode vẫn tách khỏi goal interface
để người vận hành giữ quyền kiểm soát chuyến bay.

\section{Cost MPPI dùng trong thí nghiệm}

MPPI lấy mẫu nhiều chuỗi lệnh, tính tổng cost $J_i$ cho rollout $i$, rồi ưu tiên
rollout có cost thấp. Trọng số được viết gọn như sau \cite{williams2018itmppi}:
\begin{equation}
  a_i = \exp\!\left(-\frac{J_i-J_{min}}{\lambda}\right),
  \qquad w_i = \frac{a_i}{\sum_j a_j}.
\end{equation}
Nominal control được cập nhật bằng trung bình có trọng số của nhiễu sampled và
chỉ lệnh đầu tiên được gửi đi; quá trình lặp lại ở chu kỳ sau.

Profile bám path sử dụng tổng cost đơn giản
\begin{equation}
  J = J_{path}+J_{vel}+J_{input}+J_{smooth}+J_{collision}.
\end{equation}
Ý nghĩa từng thành phần:
\begin{align}
J_{path} &= c_p\sum_k \|p_k-p_k^{ref}\|^2,
&\text{bám vị trí trên global path},\\
J_{vel} &= c_v\sum_k \|v_k-v_k^{ref}\|^2,
&\text{bám tốc độ tham chiếu},\\
J_{input} &= \sum_k u_k^T R u_k,
&\text{tránh lệnh điều khiển quá lớn},\\
J_{smooth} &= \sum_k (u_k-u_{k-1})^T R_{\Delta}(u_k-u_{k-1}),
&\text{giảm thay đổi lệnh},\\
J_{collision} &= C\sum_k \mathbf{1}[d_k<r_{safe}],
&\text{phạt rollout đi vào vùng va chạm}.
\end{align}

Trong đó $p_k^{ref},v_k^{ref}$ là reference tại đúng bước dự đoán; $d_k$ là
khoảng cách từ rollout tới point cloud; $C$ là hệ số phạt lớn. Cấu trúc
reference/input/collision được đối chiếu với Minařík et al.
\cite{minarik2024agilemppi}. Việc ánh xạ sang velocity command, bán kính point-cloud,
weights số và conditioner là lựa chọn engineering của project, không phải bản
sao nguyên tham số paper.

PA-MPPI v0 bổ sung map ba trạng thái $\{unknown,free,occupied\}$. Khi goal bị
che khuất, perception term ưu tiên trajectory mở rộng vùng nhìn thấy; collision
term coi vùng chưa xác nhận free là không an toàn. Ý tưởng này theo PA-MPPI
\cite{zhai2026pamppi}, nhưng mapper NumPy/Torch hiện tại chỉ là prototype và
không được gọi là reproduction đầy đủ.

\section{Thử nghiệm Gazebo/SITL}

Ba thí nghiệm dùng cùng velocity-level controller, seed 7, Gazebo Harmonic và
ArduPilot commit \code{f808f78c}. Global path được cung cấp trước; LiDAR và
odometry đi qua pipeline live. Bảng chỉ báo cáo log đã ghi, không nội suy sang
profile demo mới.

\begin{table}[h]
\centering\small
\caption{Kết quả live đã ghi với profile paper-cost ngày 13/09/2026.}
\begin{tabularx}{\textwidth}{L{0.20\textwidth}rrrrY}
\toprule
\textbf{Map} & \textbf{Cycles} & \textbf{Final error} & \textbf{Path} & \textbf{Min clearance} & \textbf{Compute mean/p95}\\
\midrule
Right angle & 128 & 0.235 m & 23.15 m & 2.19 m & 14.5/14.9 ms\\
Narrow gate & 192 & 0.230 m & 35.50 m & 2.07 m & 13.3/14.7 ms\\
Slalom & 193 & 0.220 m & 37.67 m & 1.30 m & 13.6/15.5 ms\\
\bottomrule
\end{tabularx}
\end{table}

Không run nào ghi deadline miss, \code{hold-brake} hoặc \code{recover-brake}.
Tuy nhiên, clearance là khoảng cách từ tâm odometry tới điểm LiDAR gần nhất,
chưa trừ footprint UAV. Giá trị 1.30 m ở slalom thấp hơn collision radius cấu
hình 1.50 m; do đó cần kiểm tra sensor latency, footprint và nhiều seed trước
khi kết luận safety.

Trong map narrow-gate, route demo đi qua cửa tường lớn rồi vòng phía bắc hai
cột; nó không xuyên khe nhỏ giữa hai cột. Route xuyên khe nhỏ là stress-test,
từng gây local minimum hoặc hard brake, nên không được trình bày như kết quả
demo thành công.

\section{Tune cho demo và bài học}

Profile \code{mppi\_demo\_smooth.yaml} giữ nguyên năm nhóm cost trên nhưng đổi
các tham số interface: reference speed $1.0\rightarrow0.85$ m/s,
max acceleration $1.5\rightarrow0.9$ m/s$^2$, command smoothing
$0.45\rightarrow0.35$, waypoint radius $1.0\rightarrow1.25$ m và goal radius
$0.25\rightarrow0.35$ m. Sweep candidate offline cho thấy RMS gia tốc lệnh
giảm khoảng 18\% trên các slalom seed đã chạy. Đây là evidence offline; chưa có
log live để gán kết quả Gazebo cho profile mới.

Các bài học thực nghiệm:
\begin{enumerate}
  \item kiểm tra bề rộng khả thi sau khi trừ footprint và uncertainty trước
  khi tune weights;
  \item global path quyết định homotopy, còn MPPI tối ưu chuyển động cục bộ;
  \item tăng samples/horizon không tự phá được local minimum hình học;
  \item terminal gate quá chặt làm UAV đứng sát goal nhưng node chưa kết thúc;
  \item mọi thay đổi YAML yêu cầu restart planner; mỗi run phải dùng log riêng;
  \item một video hoặc một seed chỉ chứng minh pipeline chạy được, không chứng
  minh robustness.
\end{enumerate}

\section{Kết luận và công việc tiếp theo}

Tuần này đã hoàn thiện luồng velocity-level từ sensing đến MAVLink, bổ sung
global-path cost, visualization sampled trajectories và goal interface trên
RViz. Ba fixed-route run cho thấy controller đạt terminal trong SITL; unit test
hiện đạt 41/41. Phần còn thiếu là live verification cho profile mượt và RViz
click, ablation MPPI/PA-MPPI trong môi trường unknown, rigid-body hover gate và
đánh giá trên edge computer.

Thứ tự tiếp theo là: (1) ghi log RViz-click trên ba map; (2) chạy nhiều seed với
cùng route và safety limits; (3) đo success, clearance, smoothness và compute
p95; (4) chỉ sau đó mới chuyển sang rigid-body và hardware gate.

\bibliographystyle{unsrt}
\bibliography{pa_mppi_sources}

\end{document}
